% Generated draft: requires mathematical and source review.
\documentclass[11pt]{article}
\usepackage[margin=25mm]{geometry}
\usepackage{fontspec}
\setmainfont{Latin Modern Roman}
\usepackage{amsmath,amssymb,mathtools,unicode-math}
\setmathfont{Latin Modern Math}
\usepackage{enumitem,needspace,xurl,hyperref}
\setlength{\parindent}{0pt}
\setlength{\parskip}{5pt}
\setlength{\emergencystretch}{3em}
\newcommand{\N}{\mathbb{N}}
\newcommand{\Z}{\mathbb{Z}}
\newcommand{\Q}{\mathbb{Q}}
\newcommand{\R}{\mathbb{R}}
\newcommand{\C}{\mathbb{C}}
\newcommand{\F}{\mathbb{F}}
\newcommand{\A}{\mathbb{A}}
\newcommand{\PP}{\mathbb{P}}
\newcommand{\E}{\mathbb{E}}
\newcommand{\eps}{\varepsilon}
\newcommand{\dd}{\,\mathrm{d}}
\begin{document}
\section*{UnsolvedMath ID: 30003094}
\textbf{Source problem code:} catalogue entry
\subsection{Definitions and mathematical statement}
The supplied problem asks for an extension of the error analysis of the dynamical low-rank splitting scheme to spatial discretizations of partial differential equations.

For notation only, let \(h\) index a spatial discretization. If \(u_h(t)\) denotes a reference evolution and \(Y_h^n\) denotes the low-rank splitting approximation at time \(t_n\), an error quantity could be written as
\[
 e_h^n=\lVert u_h(t_n)-Y_h^n\rVert_h,
\]
provided that both states belong to a common space equipped with a specified norm \(\lVert\cdot\rVert_h\). This notation does not specify a PDE class, a discretization, a rank constraint, a splitting algorithm, or a reference solution; these are absent from the supplied record.

\textbf{Requested mathematical task.} Establish an error analysis for the intended dynamical low-rank splitting scheme when applied to spatially discretized PDEs, addressing the obstruction described in the supplied curation: the existing estimate depends on a Lipschitz constant and does not apply to these spatial discretizations.

\textbf{Unresolved qualifications.} The record supplies neither the existing estimate nor its hypotheses. It does not specify whether the desired error is relative to the semidiscrete PDE evolution, an exact low-rank evolution, or the original PDE solution. Nor does it specify an error rate, a time interval, admissible initial or boundary conditions, or the required dependence on spatial resolution and rank. Thus no fully quantified error theorem can be extracted without additional source information. In particular, a bound uniform in spatial resolution is a possible interpretation of the intended extension, not an explicit requirement stated in the record.
\subsection{Short English statement}
Develop an error analysis for the intended dynamical low-rank splitting method applied to spatially discretized PDEs, overcoming the reported limitation involving a Lipschitz constant. The PDE class, numerical scheme, assumptions, error norm, and desired estimate remain unspecified, so the supplied record determines a research task rather than a fully quantified conjecture.
\subsection{Sources}
\begin{itemize}
\item Local UnsolvedMath registry: \texttt{data/unsolvedmath/problems.json}, permanent ID \(30003094\), source problem code \texttt{catalogue entry}, title ``Low-Rank Splitting Error Analysis for Discretized PDEs.'' The supplied record belongs to the Oberwolfach Reports Open Problems collection and gives proposed year 2016; no individual proposer is identified.
\item Supplied report citation: \emph{Geometric Numerical Integration} (2016), Oberwolfach Reports, pp.\ 869--948, DOI 10.4171/owr/2016/18. \url{https://doi.org/10.4171/owr/2016/18}. This is also the supplied literature source labeled ``OWR 2016 report.'' No precise problem page, section, or theorem reference is supplied.
\item Literature assessment checked at \texttt{2026-08-21}: the supplied curation labels the problem open, reports that the existing estimate depends on a Lipschitz constant and fails for spatial discretizations of PDEs, and states that no authoritative later result supplying the missing general error analysis was located. This is supplied dated curation, not an independently verified literature assessment. The background explicitly requests independent verification of time-sensitive status and cited claims.
\item The registry labels the statement \texttt{corrected\_verified} and reports a manual semantic, typographic, and LaTeX review. These are supplied editorial qualifications; they do not resolve the mathematical specifications missing from this record.
\end{itemize}
\end{document}
